\chapter{附着于新形式的 Galois 表示}

% ===== PDF p.070 =====
\noindent\textbf{下一个目标}：
\begin{itemize}
\item 在上同调中实现新形式。
\item 完成「附着于权 $k=2$ 的新形式的 Galois 表示」这一构造。
\item 对 $k\ge2$ 的情形给出如何继续的提示。
\end{itemize}

我们将借助微分形式，把（权 $2$ 的）模形式对应到奇异上同调中的类。为此需要回忆 de Rham 上同调的一些事实。

\noindent\textbf{de Rham 上同调}（参见 \texttt{[BT13]}）：
\begin{itemize}
\item 设 $X$ 为实流形。
\item 令 $\C^\infty(X)$ 为光滑函数 $f\colon X\to\C$ 构成的空间。
\item 令 $A^i(X)$ 为 $X$ 上光滑 $i$-形式构成的空间，$i\ge0$。
\item $X$ 的 \textbf{de Rham 复形}是复形
\[
A^\bullet(X)\colon\quad A^0(X)=\C^\infty(X)\xrightarrow{\ \mathrm{d}\ }A^1(X)\xrightarrow{\ \mathrm{d}\ }A^2(X)\to\cdots
\]
其中 $\mathrm{d}$ 为外微分。
\item $X$ 的 de Rham 上同调则定义为 de Rham 复形的上同调
\[
H^*_{\mathrm{dR}}(X):=H^*(A^\bullet(X)).
\]
\item 引人注目的是，de Rham 比较同构蕴涵
\[
H^*_{\mathrm{dR}}(X)\cong H^*(X,\C),
\]
右边是 $X$ 上常层 $\C$ 的层上同调。这有个有趣的推论：de Rham 上同调是 $X$ 的纯拓扑不变量。
\begin{itemize}
\item 证明略述：把开集 $U\subseteq X$ 送到 $A^\bullet(U)$，给出 $X$ 上的层复形 $\mathcal A^\bullet_X$，它是常层 $\C$ 的一个消解（由 Poincaré 引理）。用单位分解的论证（需要 $X$ 仿紧）可知，层 $\mathcal A^k_X$（$k\ge0$）都是 flasque（松弛）层。于是 $R\Gamma(X,\C)\cong\Gamma(X,\mathcal A^\bullet_X)$。
\end{itemize}
\item 此外，
\[
H^*(X,\C)\cong H^*_{\mathrm{sing}}(X,\C),
\]
右边是 $X$ 的以 $\C$ 为系数的奇异上同调。
\item 若离散群 $\Gamma$ 真不连续且自由地作用在 $X$ 上，则
\[
A^\bullet(\Gamma\backslash X)\cong A^\bullet(X)^\Gamma=\bigl(A^0(X)^\Gamma=\C^\infty(X)^\Gamma\xrightarrow{\ \mathrm{d}\ }A^1(X)^\Gamma\xrightarrow{\ \mathrm{d}\ }A^2(X)^\Gamma\to\cdots\bigr),
\]
也就是说，de Rham 复形中的不变量计算商 $\Gamma\backslash X$ 的上同调。注意：由于群上同调可能出现，这一陈述若放到上同调群层面便不成立。

% ===== PDF p.071 =====
\item 最后，我们令
\[
H^*_{\mathrm{dR}}\bigl(\GL_2(\Q)\backslash(\GL_2(\Af)\times\HH^\pm)\bigr):=\varinjlim_{K}H^*_{\mathrm{dR}}\bigl(\GL_2(\Q)\backslash(\GL_2(\Af)/K\times\HH^\pm)\bigr)
\]
以及
\[
H^*\bigl(\GL_2(\Q)\backslash(\GL_2(\Af)\times\HH^\pm),\C\bigr):=\varinjlim_{K}H^*\bigl(\GL_2(\Q)\backslash(\GL_2(\Af)/K\times\HH^\pm),\C\bigr),
\]
其中余极限取遍全部紧开子群 $K\subseteq\GL_2(\Af)$，过渡映射由拉回诱导。
\item 于是对 $i\ge0$，取余极限即得（$\GL_2(\Af)$-等变的）de Rham 比较
\[
H^i_{\mathrm{dR}}\bigl(\GL_2(\Q)\backslash(\GL_2(\Af)\times\HH^\pm)\bigr)\cong H^i\bigl(\GL_2(\Q)\backslash(\GL_2(\Af)\times\HH^\pm),\C\bigr).
\]
\znote{这里取余极限得到的大空间就是第 8 讲的 $\widehat X(\C)=\GL_2(\Q)\backslash(\GL_2(\Af)\times\HH^\pm)$，所以上面两条定义左边正是概形 $\widehat X$ 的上同调。下文把 $\C$ 换成 $\Q_\ell$、把 de Rham 上同调换成平展上同调，正是为了让 Galois 群在它上面作用。[来源：本讲义第 8 讲（PDF p.063）对 $\widehat X(\C)$ 的刻画。]}
\end{itemize}

\noindent\textbf{要点}：
\begin{itemize}
\item 可以用（闭的）微分形式构造 $H^*(X,\C)$ 中的类。
\end{itemize}

\noindent\textbf{附着于权 2 模形式的上同调类}：
\begin{itemize}
\item 注意
\[
\GL_2(\Q)\backslash(\GL_2(\Af)\times\HH^\pm)\cong\GL_2(\Q)^+\backslash(\GL_2(\Af)\times\HH),
\]
其中 $\GL_2(\Q)^+\subseteq\GL_2(\Q)$ 是行列式为正的元素构成的子群。
\item 取 $f\in H^0\bigl(\GL_2(\Q)^+\backslash(\GL_2(\Af)\times\HH),\omega^{\otimes k}\bigr)$，$k\in\Z$。
\item 我们把 $f$ 看作函数
\[
f\colon\GL_2(\Af)\times\HH\to\C,\qquad (g,z)\mapsto f(g,z),
\]
它满足
\[
f(\gamma g,\gamma z)=(cz+d)^kf(g,z),\qquad \gamma=\begin{pmatrix}a&b\\c&d\end{pmatrix}\in\GL_2(\Q)^+ .
\]
\item 假设 $f$ 的权为 $k=2$。
\item 此时 $\GL_2(\Af)\times\HH$ 上的微分形式 $f(g,z)\,\mathrm{d}z$ 是闭的，并且对 $\gamma\in\GL_2(\Q)^+$ 满足
\[
\gamma^*\bigl(f(g,z)\,\mathrm{d}z\bigr)=f(\gamma g,\gamma z)\,\gamma^*\mathrm{d}z=\det(\gamma)\,f(g,z)\,\mathrm{d}z .
\]
\item 确实：
\begin{itemize}
\item 闭性由全纯性得到，因为
\[
\mathrm{d}\bigl(f(g,z)\,\mathrm{d}z\bigr)=\frac{\partial f}{\partial z}\,\mathrm{d}z\wedge\mathrm{d}z-\frac{\partial f}{\partial\bar z}\,\mathrm{d}z\wedge\mathrm{d}\bar z=0 .
\]
\item 若 $\gamma=\begin{pmatrix}a&b\\c&d\end{pmatrix}\in\GL_2(\Q)$，则
\[
\gamma^*(\mathrm{d}z)=\frac{\det(\gamma)}{(cz+d)^2}\,\mathrm{d}z .
\]
\item 利用 $f$ 的权为 $2$，即得结论。
\end{itemize}

% ===== PDF p.072 =====
\item 令
\[
|\cdot|=|\cdot|_{\mathrm{ad\acute{e}lic}}\colon\Q^\times\backslash\A^\times\to\R_{>0}
\]
为 adélic 绝对值，即
\[
|(x_2,x_3,\dots,x_\infty)|:=\prod_p|x_p|_p\cdot|x_\infty|_\infty,
\]
其中 $|\cdot|_p\colon\Q_p\to\R_{\ge0}$，$x\mapsto p^{-v_p(x)}$ 是 $p$ 进绝对值，$|\cdot|_\infty\colon\R\to\R_{\ge0}$ 是实绝对值。
\item $|\cdot|$ 定义了 $\GL_2(\A)$ 的特征
\[
\chi_{no}:=|\cdot|\circ\det=|\det|_{\mathrm{ad\acute{e}lic}},
\]
它在 $\GL_2(\Q)$ 上平凡。
\item 于是对
\[
\tilde f(g,z):=\chi_{no}((g,1))f(g,z)
\]
（其中 $(g,1)\in\GL_2(\A)\cong\GL_2(\Af)\times\GL_2(\R)$），$\GL_2(\Af)\times\HH$ 上的形式
\[
\eta_f:=\tilde f(g,z)\,\mathrm{d}z
\]
是 $\GL_2(\Q)^+$-等变的。
\item 我们由此得到映射
\[
\alpha\colon M_2\longrightarrow H^1_{\mathrm{dR}}\bigl(\GL_2(\Q)^+\backslash(\GL_2(\Af)\times\HH)\bigr)\cong H^1\bigl(\GL_2(\Q)^+\backslash(\GL_2(\Af)\times\HH),\C\bigr),
\]
它把 $f$ 送到 $[\eta_f]$；这个映射不是 $\GL_2(\Af)$-等变的。
\item 对光滑 $\GL_2(\Af)$-表示 $V$ 与光滑特征 $\chi\colon\GL_2(\Af)\to\C^\times$，定义
\[
V(\chi):=V\otimes_\C\chi .
\]
\znoteplain{这里的扭是表示层面的一维扭：作为向量空间 $V(\chi)=V$，而 $g\in\GL_2(\Af)$ 的作用改为「先按 $V$ 作用、再乘 $\chi(g)$」。本讲此后反复用到这一记号，例如 $M_2(\chi_{no})$、$\pi(\chi_{no})$。}
\item 于是
\[
\alpha\colon M_2(\chi_{no})\longrightarrow H^1\bigl(\GL_2(\Q)^+\backslash(\GL_2(\Af)\times\HH),\C\bigr)
\]
是 $\GL_2(\Af)$-等变的。
\item 设 $K\subseteq\GL_2(\Af)$ 为紧开子群，并令
\[
X_K:=\GL_2(\Q)^+\backslash(\GL_2(\Af)/K\times\HH),
\]
其典范紧化为
\[
X^*_K:=\GL_2(\Q)^+\backslash(\GL_2(\Af)/K\times\HH^*),
\]
其中 $\HH^*:=\HH\cup\Pone(\Q)$（带 Satake 拓扑）。
\item 若 $f\in S_2(K)$，则 $\eta_f$ 延拓为 $X^*_K$ 上的全纯微分形式。
\begin{itemize}
\item 确实：若 $q=e^{2\pi iz}$，则
\[
\mathrm{d}q=2\pi iq\,\mathrm{d}z,\qquad\text{即}\qquad \mathrm{d}z=\frac{1}{2\pi iq}\,\mathrm{d}q .
\]
于是把 $f(g,z)$ 在每个尖点处写成 Fourier 展开，即写成 $q$ 的函数。
\end{itemize}

% ===== PDF p.073 =====
\item 下图交换：
\[
\begin{tikzcd}
H^1_{\mathrm{dR}}(X^*_K)\arrow[r,"\simeq"]\arrow[d,"\simeq"] & H^1_{\mathrm{dR}}(X_K)\arrow[d,"\simeq"]\\
H^1(X^*_K,\C)\arrow[r,"\simeq"] & H^1(X_K,\C)
\end{tikzcd}
\]
\item 此外，态射
\[
H^1_c(X_K,\C)\longrightarrow H^1(X^*_K,\C)=H^1_c(X^*_K,\C)
\]
是满射，因为其余核嵌入 $H^1(\{\text{尖点}\},\C)=0$。
\item 于是 $\alpha(S_2(K))$ 落在 $X_K$ 的内上同调
\[
\tilde H^1(X_K,\C)\cong H^1(X^*_K,\C)
\]
中；按定义，后者是紧支上同调 $H^1_c(X_K,\C)\to H^1(X_K,\C)$ 的像。
\item 注意上述讨论对反全纯模形式（即模形式的复共轭）同样适用。
\item 由此得到映射
\[
\bar\alpha\colon\bar S_2(K)\longrightarrow\tilde H^1(X_K,\C),\qquad f\mapsto[\chi_{no}f\,\mathrm{d}\bar z],
\]
其中 $\bar S_2(K)$ 表示反全纯模形式构成的 $\C$-线性空间。
\item 我们可以过渡到无穷级，得到态射
\[
\bar\alpha\oplus\alpha\colon\bar S_2\oplus S_2\longrightarrow\tilde H^1\bigl(\GL_2(\Q)\backslash(\GL_2(\Af)\times\HH^\pm),\C\bigr),
\]
其记号（希望）自明。
\end{itemize}

\noindent\textbf{定理 9.1}（Eichler--Shimura）\phantomsection\label{thm:9.1}\quad 映射
\[
\alpha\oplus\bar\alpha\colon S_2(\chi_{no})\oplus\bar S_2(\chi_{no})\longrightarrow\tilde H^1\bigl(\GL_2(\Q)\backslash(\GL_2(\Af)\times\HH^\pm),\C\bigr)
\]
是 $\GL_2(\Af)$-等变同构。

\begin{itemize}
\item 该陈述等价于它对所有（充分小的）紧开子群 $K\subseteq\GL_2(\Af)$ 的类似陈述。
\item 于是固定某个紧开子群 $K\subseteq\GL_2(\Af)$。
\item 证明将利用 $\cup$-积配对
\[
H^1_c(X_K,\C)\times H^1(X_K,\C)\longrightarrow H^2_c(X_K,\C)\xrightarrow{\ \int\ }\C,
\]
在 de Rham 比较之下，它由微分形式的配对
\[
(\eta_1,\eta_2)\mapsto\int_{X_K}\eta_1\wedge\eta_2
\]
诱导，其中 $\eta_1$ 具有紧支集（参见 \texttt{[BT13]}）。
\item 我们将用到如下观察：
\begin{itemize}
\item 紧支上同调可由 $H^*_c(X_K,\C)\cong H^*(A^\bullet_c(X_K))$ 计算，其中 $A^\bullet_c$ 表示紧支集微分形式。
\item 典范映射
\begin{equation}\tag{2}\label{eq:9.2}
A^\bullet_c(X_K)\longrightarrow A^\bullet_{\mathrm{rd}}(X_K)
\end{equation}
是拟同构，右边表示速降微分形式，参见 \texttt{[Bor80, Theorem 2]}。
\end{itemize}
\item 对 $f\in S_2$，形式 $\eta_f$ 速降。
\begin{itemize}
% ===== PDF p.074 =====
\item 利用 $|e^{2\pi inz}|=e^{-2\pi n\operatorname{Im}(z)}$ 当 $\operatorname{Im}(z)\to\infty$ 且 $n\ge1$ 时迅速衰减。
\end{itemize}
\item 于是 $\alpha$ 自然提升为映射
\[
\tilde\alpha\colon S_2(K)\longrightarrow H^1_c(X_K,\C).
\]
\item 取 $f\in S_2(K)$、$g\in\bar S_2(K)$，则：
\begin{itemize}
\item $\tilde\alpha(f)\cup\alpha(f)=0$，因为 $\mathrm{d}z\wedge\mathrm{d}z=0$。
\item 对 $g$ 同理。
\item 用 Stokes 定理可证
\[
\operatorname{tr}\bigl(\tilde\alpha(f)\cup\bar\alpha(g)\bigr)=\int_{X_K}\eta_f\wedge\eta_g .
\]
\item 由此可知，Poincaré 配对（相差一个标量）诱导出 $S_2$ 上的 Petersson 标量积，又称 $L^2$-配对 $\langle-,-\rangle$。
\end{itemize}
\item 现在可以证明 $\beta:=\alpha\oplus\bar\alpha$ 是单射。
\item 确实：
\begin{itemize}
\item 若 $\beta(f+g)=0$，则
\[
\beta(f+g)\cup\alpha(f)=\langle f,f\rangle=0,
\]
即 $f=0$。对 $g$ 同理。
\end{itemize}
\item 现在验证 $\tilde h^1(X_K,\C)=2\dim S_2(K)$。一旦成功，Eichler--Shimura 同构\idx{Eichler--Shimura 同构}的证明便完成了。
\item 即（我们用小写字母表示维数）：
\begin{itemize}
\item $\tilde h^1(X_K,\C)=h^1(X^*_K,\C)=2h^0(X^*_K,\C)-\chi_{\mathrm{top}}(X^*_K)$。
\item 这里 $\chi_{\mathrm{top}}$ 是拓扑 Euler 特征标。
\item 另一方面\znoteplain{原文脚注 15：我们隐含地使用了 $\omega^{\otimes2}(-\text{尖点})\cong\Omega^1$。原文作 $\omega^2$，按第 7 讲的记号应理解为 $\omega^{\otimes2}$。}：
\[
\begin{aligned}
\dim S_2(K)&=h^0\bigl(X^*_K,\Omega^1_{X^*_K}\bigr)\\
&\overset{\mathrm{RR}}{=}\deg\bigl(\Omega^1_{X^*_K}\bigr)+\chi_{\mathrm{hol}}(X^*_K)+h^1\bigl(X^*_K,\Omega^1_{X^*_K}\bigr)\\
&=-\chi_{\mathrm{hol}}(X^*_K)+h^0(X^*_K,\C).
\end{aligned}
\]
\znote{中间两步的来历：对曲线 $X^*_K$ 上的 $\Omega^1$ 用 Riemann--Roch 即得第一个等号后的三项；最后一步用 $\deg\Omega^1_{X^*_K}=2g-2=-2\chi_{\mathrm{hol}}(X^*_K)$ 与 Serre 对偶 $h^1(X^*_K,\Omega^1)=h^0(X^*_K,\OO)=1$，于是 $\dim S_2(K)=1-\chi_{\mathrm{hol}}(X^*_K)=-\chi_{\mathrm{hol}}(X^*_K)+h^0(X^*_K,\C)$。[来源：R.~Hartshorne, \emph{Algebraic Geometry}, GTM 52, IV.1.3（Riemann--Roch 定理）与 III.7.6（Serre 对偶定理）。]}
\item 这里 $\chi_{\mathrm{hol}}$ 是全纯 Euler 特征标。
\item $\chi_{\mathrm{top}}(X^*_K)=2\chi_{\mathrm{hol}}(X^*_K)$。
\end{itemize}
\end{itemize}

至此，我们完成了为新形式附着 Galois 表示的既定路线（至少在权 $2$ 的情形）。

\noindent\textbf{附着于新形式的 Galois 表示}：
\begin{itemize}
\item 回忆：前两讲我们引入了概形
\[
\widehat X\to\Spec(\Q),
\]
它带有 $\GL_2(\Af)$-作用，使得自然地
\[
\widehat X(\C)\cong\GL_2(\Q)\backslash(\GL_2(\Af)\times\HH^\pm).
\]
\item 对紧开（且充分小）的 $K\subseteq\GL_2(\Af)$，得到拟射影光滑曲线
\[
X_K\to\Spec(\Q),\qquad X_K(\C)\cong\GL_2(\Q)\backslash(\GL_2(\Af)/K\times\HH^\pm).
\]
% ===== PDF p.075 =====
（注意与上一节记号的冲突：那里 $X_K$ 就等于右边的商。）\znoteplain{第 8 讲的 $X_K$ 是 $\Spec(\Q)$ 上的曲线 $\widehat X/K$（概形），而本讲从第 72 页起 $X_K$ 指的是它的 $\C$-点集 $\GL_2(\Q)^+\backslash(\GL_2(\Af)/K\times\HH)$；两者是同一对象的两种看法，只是记法相撞。译文保留原记法：凡讲「曲线、概形」时指前者，凡讲「复流形、上同调」时指后者。}
\item 固定一个素数 $\ell$。
\item 对 $K\subseteq\GL_2(\Af)$，可以考虑内平展上同调
\[
\tilde H^1_{\mathrm{\acute{e}t}}\bigl(X_{K,\Q},\Q_\ell\bigr):=\operatorname{Im}\Bigl(H^1_{c,\mathrm{\acute{e}t}}\bigl(X_{K,\Q},\Q_\ell\bigr)\to H^1_{\mathrm{\acute{e}t}}\bigl(X_{K,\Q},\Q_\ell\bigr)\Bigr)\cong H^1_{\mathrm{\acute{e}t}}\bigl(X^*_{K,\Q},\Q_\ell\bigr).
\]
\item 取极限，得到大的上同调空间
\[
\tilde H^1_{\mathrm{\acute{e}t}}\bigl(\widehat X_{\Q},\Q_\ell\bigr):=\varinjlim_{K}\tilde H^1_{\mathrm{\acute{e}t}}\bigl(X_{K,\Q},\Q_\ell\bigr).
\]
\item 由 $H^1_{c,\mathrm{\acute{e}t}}$、$H^1_{\mathrm{\acute{e}t}}$ 的自然性，这个空间带有
\[
\Gal(\overline\Q/\Q)\times\GL_2(\Af)
\]
的作用，即 $\Gal(\overline\Q/\Q)$ 的一个作用与 $\GL_2(\Af)$ 的一个作用，且这两个作用交换。
\item 固定一个同构 $\iota\colon\Ql\xrightarrow{\ \sim\ }\C$。
\item 利用平展比较定理（以及 $\iota$），得到同构
\[
\tilde H^1_{\mathrm{\acute{e}t}}\bigl(\widehat X_{\Q},\Q_\ell\bigr)\cong H^1_{\mathrm{\acute{e}t}}\bigl(\widehat X_{\C},\Q_\ell\bigr)\cong\tilde H^1\bigl(\widehat X(\C),\Q_\ell\bigr)\otimes_{\Q_\ell,\iota}\C\cong\tilde H^1\bigl(\widehat X(\C),\C\bigr).
\]
\item 这些同构都是 $\GL_2(\Af)$-等变的。
\item 设 $\pi\subseteq S_2$ 为不可约 $\GL_2(\Af)$-表示。
\item 借助 $\iota$，我们把每个 $\C$ 上的光滑 $\GL_2(\Af)$-表示都看作 $\Ql$ 上的光滑表示。
\item 空间
\[
\tilde\rho_\pi:=\Hom_{\GL_2(\Af)}\bigl(\pi(\chi_{no}),\tilde H^1_{\mathrm{\acute{e}t}}(\widehat X_{\Q},\Q_\ell)\bigr)
\]
是二维的：这由 Eichler--Shimura 同构
\[
\tilde H^1(\widehat X(\C),\C)\cong S_2(\chi_{no})\oplus\bar S_2(\chi_{no})
\]
以及 $S_2$ 作为 $\GL_2(\Af)$-表示的重数一性质得到。
\item 总之，作为 $\Gal(\overline\Q/\Q)\times\GL_2(\Af)$-表示，
\begin{equation}\tag{3}\label{eq:9.3}
\tilde H^1_{\mathrm{\acute{e}t}}\bigl(\widehat X_{\Q},\Q_\ell\bigr)\cong\bigoplus_\pi\tilde\rho_\pi\otimes_{\Q_\ell}\pi(\chi_{no}),
\end{equation}
其中求和取遍全部不可约 $\GL_2(\Af)$-表示 $\pi\subseteq S_2$（注意从右边到左边有自然的态射）。
\end{itemize}

前面已试图说明：新形式的 Langlands 互反之所以重要，在于 Hecke 特征值与 Frobenius 的迹相合。下面就来讨论这一点。

\noindent\textbf{平展上同调中的 Eichler--Shimura 关系}：
\begin{quote}
我们需要说明：Hecke 特征值与 Frobenius 的迹相合。
\end{quote}

\begin{itemize}
\item 回忆：对 $\pi\subseteq S_k$ 及其 Hecke 特征系统
\[
\bigl\{a_p(\pi),\ b_p(\pi)=\chi(p)p^{k-1}\bigr\}_{p\notin S},
\]
我们想附着二维 $\ell$-进表示
\[
\rho_\pi\colon\Gal(\overline\Q/\Q)\longrightarrow\GL_2(\Q_\ell),
\]
使得对每个 $p\notin S$，$p$ 处的算术 Frobenius $\Frob_p^{\mathrm{arith}}$ 的特征多项式为
\[
X^2-a_p(\pi)X+\chi(p)p^{k-1}
\]
% ===== PDF p.076 =====
（以下我们省略 $\iota$ 与下标 $\pi$），即
\[
\bigl(\Frob_p^{\mathrm{arith}}\bigr)^2-a_p(\pi)\Frob_p^{\mathrm{arith}}+\chi(p)p^{k-1}=0 .
\]
\item 在（式 (3)）中取 $K_1(N)$-不变量，得到（对 $N\ge3$）
\[
H^1_{\mathrm{\acute{e}t}}\bigl(X^*_{\Gamma_1(N),\Q},\Q_\ell\bigr)\cong\bigoplus_\pi\tilde\rho_\pi\otimes_{\Q_\ell}\pi(\chi_{no})^{K_1(N)}
\]
（注意 $X_{K_1(N)}\cong X_{\Gamma_1(N)}$）。
\item 设 $p$ 为素数，$p\nmid\ell N$。
\item 固定 $\overline\Q$ 中位于 $p$ 之上的一个位。这确定了 $\F_p$ 的一个代数闭包 $\overline\F_p$。
\item 回忆：对域 $L$ 上的 abelian 簇 $A$ 与素数 $\ell$，记
\[
T_\ell A=\varprojlim_n A[\ell^n](\overline L),\qquad V_\ell A=T_\ell A\otimes_{\Z_\ell}\Q_\ell
\]
分别为 $\ell$-进 Tate 模与有理化的 $\ell$-进 Tate 模。
\end{itemize}

\noindent\textbf{引理 9.2}\phantomsection\label{lem:9.2}\quad 设 $p$ 为素数，$p\nmid\ell N$。则
\[
H^1_{\mathrm{\acute{e}t}}\bigl(X^*_{\Gamma_1(N),\Q},\Q_\ell\bigr)=V_\ell\Pic^0\bigl(X^*_{\Gamma_1(N),\Fq}\bigr)^\vee,
\]
其中 $(-)^\vee$ 表示 $\Q_\ell$-对偶。

\noindent\textbf{证明}\quad Kummer 序列 $1\to\mu_{\ell^n}\to\Gm\xrightarrow{\ell^n}\Gm\to1$ 蕴涵
\[
H^1_{\mathrm{\acute{e}t}}\bigl(X^*_{\Gamma_1(N),\overline\Q},\Q_\ell(1)\bigr)\cong V_\ell\Pic^0\bigl(X^*_{\Gamma_1(N),\overline\Q}\bigr).
\]
Weil 配对给出典范同构
\[
V_\ell\Pic^0\bigl(X^*_{\Gamma_1(N),\overline\Q}\bigr)^\vee\cong V_\ell\Pic^0\bigl(X^*_{\Gamma_1(N),\overline\Q}\bigr)(-1).
\]
最后，由于 $\ell\neq p$，
\[
V_\ell\Pic^0\bigl(X^*_{\Gamma_1(N),\overline\Q}\bigr)\cong V_\ell\Pic^0\bigl(X^*_{\Gamma_1(N),\overline\Q_p}\bigr)\cong V_\ell\Pic^0\bigl(X^*_{\Gamma_1(N),\Fq}\bigr),
\]
因为挠点可以提升。\hfill$\square$

\begin{itemize}
\item 特别地，Galois 表示
\[
H^1_{\mathrm{\acute{e}t}}\bigl(X^*_{\Gamma_1(N),\Q},\Q_\ell\bigr)
\]
在 $p$ 处非分歧。
\item 在 $V_\ell\Pic^0(X^*_{\Gamma_1(N),\Fq})$ 上有 Eichler--Shimura 关系，回忆它写作
\[
F^2-\tilde T_pF+p\tilde S_p=0,
\]
其中 $F$ 是算术 Frobenius。\znoteplain{为什么在 $\Pic^0$ 的 Tate 模上关系对算术 Frobenius 成立，而命题 9.3 在 $H^1_{\mathrm{\acute{e}t}}$ 上却对几何 Frobenius 成立？由引理 9.2，$H^1_{\mathrm{\acute{e}t}}$ 是 $V_\ell\Pic^0$ 的 $\Q_\ell$-对偶，对偶把算术 Frobenius 与几何 Frobenius 互换；而转置不改变特征多项式，故两处给出的是同一个多项式。}
\end{itemize}

\noindent\textbf{命题 9.3}\phantomsection\label{prop:9.3}\quad 设 $p$ 为素数，$p\nmid\ell N$。则在 $H^1_{\mathrm{\acute{e}t}}(X^*_{\Gamma_1(N),\Q},\Q_\ell)$ 上有
\[
\bigl(\Frob_p^{\mathrm{geom}}\bigr)^2-\tilde T_p\Frob_p^{\mathrm{geom}}+p\tilde S_p=0 .
\]

\begin{itemize}
\item 对归一化新形式 $f\in S_2(\Gamma_0(N),\chi)$，我们最终定义
\[
\rho_f:=(\tilde\rho_\pi)^\vee,
\]
其中 $\pi$ 是由 $f$ 生成的不可约 $\GL_2(\Af)$-表示，而
\[
\tilde\rho_\pi:=\Hom_{\GL_2(\Af)}\bigl(\pi(\chi_{no}),\tilde H^1_{\mathrm{\acute{e}t}}(\widehat X_{\Q},\Q_\ell)\bigr).
\]
\item 关于 $\rho_f$ 的一些分析，参见 \texttt{[Rib77]}：
\begin{itemize}
% ===== PDF p.077 =====
\item 把（归一化）新形式 $f\in S_2(\Gamma_0(N),\chi)$ 写成 Fourier 展开
\[
f(q)=\sum_{n\ge1}a_nq^n .
\]
则对 $p\nmid\ell N$，$\Frob_p^{\mathrm{arith}}$ 在 $\rho_f$ 上的特征多项式为
\[
X^2-a_pX+\chi(p)p .
\]
\item 确实：
\begin{itemize}
\item 由 $\Frob_p^{\mathrm{geom}}$ 的 Eichler--Shimura 关系（的对偶），在 $\rho_f$ 上得到
\[
\bigl(\Frob_p^{\mathrm{arith}}\bigr)^2-\tilde T_p\Frob_p^{\mathrm{arith}}+p\tilde S_p=0 .
\]
\item 回忆 $\tilde T_p$ 在 $\pi^{K_1(N)}$ 上的特征值是 $pa_p$，而 $\tilde S_p$ 的特征值是 $\chi(p)p^2$。
\item 我们曾把 $S_2$ 上的 $\GL_2(\Af)$-作用按 $\chi_{no}=|\det(-)|$ 扭过。因此 $\tilde T_p$ 在
\[
\tilde\rho_\pi\otimes_{\Q_\ell}\pi(\chi_{no})^{K_1(N)}\subseteq H^1_{\mathrm{\acute{e}t}}\bigl(X^*_{\Gamma_1(N),\Q},\Q_\ell\bigr)
\]
上的特征值是 $a_p$，而 $\tilde S_p$ 的特征值是
\[
\frac{1}{p^2}\chi(p)p^2=\chi(p).
\]
\end{itemize}
\item $\rho_f$ 的行列式为
\[
\psi\cdot\chi_{\mathrm{cyc}},
\]
其中
\[
\chi_{\mathrm{cyc}}\colon G_\Q\to\Q_\ell^\times
\]
表示分圆特征，而
\[
\psi\colon\Q^\times\R_{>0}\backslash\A^\times\to\Q_\ell^\times\cong_{\iota}\C^\times
\]
是由 $\chi$ 决定的（有限阶）adélic 特征。
\item 确实：
\begin{itemize}
\item 用 Chebotarev 验证：对 $p\nmid\ell N$，两边在算术 Frobenius 元上一致。
\end{itemize}
\item 特别地，$\rho_f$ 是奇的，即每个复共轭的行列式为 $-1$。
\begin{itemize}
\item 这用到：对一切 $k\ge1$，$S_k(\Gamma_0(N),\chi)$ 非平凡蕴涵 $\chi(-1)=(-1)^k$。
\end{itemize}
\item $\rho_f$ 不可约，参见 \texttt{[Rib77, Theorem (2.3)]}。
\item 由于 $\rho_f$ 实现为一个 $\Q$ 上 proper、光滑概形的平展上同调，$\ell$-进 Hodge 理论中的 de Rham 比较定理蕴涵 $\rho_f|_{G_{\Q_\ell}}$ 是 de Rham 表示。
\end{itemize}
\end{itemize}

\noindent\textbf{权 $\ge2$ 的构造概略}：
\begin{itemize}
\item 不用 $\tilde H^1_{\mathrm{\acute{e}t}}(\widehat X_{\Q},\Q_\ell)$，而用
\[
\tilde H^1_{\mathrm{\acute{e}t}}(\widehat X_{\Q},\cL_k),
\]
其中
\[
\cL_k:=\Sym^{k-1}R^1f_*(\Q_\ell)
\]
是局部系统\idx{局部系统}，而
% ===== PDF p.078 =====
$f\colon E\to\widehat X$ 为万有椭圆曲线。\znote{这正是第 6 讲（PDF p.053）预告的「把常数层 $\Q_\ell$ 换成某个局部系统」：$R^1f_*\Q_\ell$ 是秩 $2$ 的局部系统，其 $(k-1)$ 次对称幂 $\cL_k$ 秩为 $k$，于是 $\tilde H^1_{\mathrm{\acute{e}t}}(\widehat X_{\Q},\cL_k)$ 里装的正是权 $k$ 的新形式。[来源：本讲义第 6 讲（PDF p.053）的提示，以及 $R^1f_*\Q_\ell$ 作为椭圆曲线 Tate 模局部系统的标准描述。]}
\item 问题：$\cL_k$ 不能延拓为紧化上的局部系统。因此需要额外的论证，才能为基变换到 $\overline\F_p$ 辩护。
\item 用类似的策略可以得到 Eichler--Shimura 关系与同构的类比。
\item 并把 de Rham 比较、平展比较定理换为带系数的版本。
\item 更多细节见 \texttt{[Del71b]}。
\end{itemize}
